\documentclass[10pt,a4paper]{amsart}
\usepackage{amsmath,amssymb,amsthm,graphicx}
\usepackage{verbatim}
\newenvironment{reference}{\comment}{\endcomment}
%\newenvironment{reference}{}{}
\newenvironment{slogan}{\comment}{\endcomment}
\newenvironment{history}{\comment}{\endcomment}

% For commutative diagrams we use Xy-pic
\usepackage[all]{xy}

% We use 2cell for 2-commutative diagrams.
\xyoption{2cell}
% Theorem environments.
%
\theoremstyle{plain}
\newtheorem{theorem}[subsection]{Theorem}
\newtheorem{proposition}[subsection]{Proposition}
\newtheorem{lemma}[subsection]{Lemma}

\theoremstyle{definition}
\newtheorem{definition}[subsection]{Definition}
\newtheorem{example}[subsection]{Example}
\newtheorem{exercise}[subsection]{Exercise}
\newtheorem{situation}[subsection]{Situation}

\theoremstyle{remark}
\newtheorem{remark}[subsection]{Remark}
\newtheorem{remarks}[subsection]{Remarks}

\numberwithin{equation}{subsection}

% Macros
%
\def\lim{\mathop{\mathrm{lim}}\nolimits}
\def\colim{\mathop{\mathrm{colim}}\nolimits}
\def\Spec{\mathop{\mathrm{Spec}}}
\def\Hom{\mathop{\mathrm{Hom}}\nolimits}
\def\Ext{\mathop{\mathrm{Ext}}\nolimits}
\def\SheafHom{\mathop{\mathcal{H}\!\mathit{om}}\nolimits}
\def\SheafExt{\mathop{\mathcal{E}\!\mathit{xt}}\nolimits}
\def\Sch{\mathit{Sch}}
\def\Mor{\mathop{\mathrm{Mor}}\nolimits}
\def\Ob{\mathop{\mathrm{Ob}}\nolimits}
\def\Sh{\mathop{\mathit{Sh}}\nolimits}
\def\NL{\mathop{N\!L}\nolimits}
\def\CH{\mathop{\mathrm{CH}}\nolimits}
\def\proetale{{pro\text{-}\acute{e}tale}}
\def\etale{{\acute{e}tale}}
\def\QCoh{\mathit{QCoh}}
\def\Ker{\mathop{\mathrm{Ker}}}
\def\Im{\mathop{\mathrm{Im}}}
\def\Coker{\mathop{\mathrm{Coker}}}
\def\Coim{\mathop{\mathrm{Coim}}}

% Boxtimes
%
\DeclareMathSymbol{\boxtimes}{\mathbin}{AMSa}{"02}

%
% Macros for moduli stacks/spaces
%
\def\QCohstack{\mathcal{QC}\!\mathit{oh}}
\def\Cohstack{\mathcal{C}\!\mathit{oh}}
\def\Spacesstack{\mathcal{S}\!\mathit{paces}}
\def\Quotfunctor{\mathrm{Quot}}
\def\Hilbfunctor{\mathrm{Hilb}}
\def\Curvesstack{\mathcal{C}\!\mathit{urves}}
\def\Polarizedstack{\mathcal{P}\!\mathit{olarized}}
\def\Complexesstack{\mathcal{C}\!\mathit{omplexes}}
% \Pic is the operator that assigns to X its picard group, usage \Pic(X)
% \Picardstack_{X/B} denotes the Picard stack of X over B
% \Picardfunctor_{X/B} denotes the Picard functor of X over B
\def\Pic{\mathop{\mathrm{Pic}}\nolimits}
\def\Picardstack{\mathcal{P}\!\mathit{ic}}
\def\Picardfunctor{\mathrm{Pic}}
\def\Deformationcategory{\mathcal{D}\!\mathit{ef}}
\usepackage{cleveref}
\crefname{originalreference}{original reference}{original references}
\Crefname{originalreference}{Original reference}{Original references}
\makeatletter
\newlabel{lemma-derived-bc}{{09LS}{source}{}{sourcepointer-1}{}}
\newlabel{lemma-derived-bc@cref}{{[originalreference][0][]09LS}{[1][0][]source}{}{}{}}
\makeatother
\crefname{originallemma}{original lemma}{original lemmas}
\Crefname{originallemma}{Original lemma}{Original lemmas}
\makeatletter
\newlabel{more-algebra-lemma-K-flat-two-out-of-three}{{06Y2}{source}{}{sourcepointer-2}{}}
\newlabel{more-algebra-lemma-K-flat-two-out-of-three@cref}{{[originallemma][0][]06Y2}{[1][0][]source}{}{}{}}
\makeatother
\crefname{originallemma}{original lemma}{original lemmas}
\Crefname{originallemma}{Original lemma}{Original lemmas}
\makeatletter
\newlabel{more-algebra-lemma-colimit-K-flat}{{06Y3}{source}{}{sourcepointer-3}{}}
\newlabel{more-algebra-lemma-colimit-K-flat@cref}{{[originallemma][0][]06Y3}{[1][0][]source}{}{}{}}
\makeatother
\crefname{originalsection}{original section}{original sections}
\Crefname{originalsection}{Original section}{Original sections}
\makeatletter
\newlabel{more-algebra-section-derived-tensor-product}{{06XY}{source}{}{sourcepointer-4}{}}
\newlabel{more-algebra-section-derived-tensor-product@cref}{{[originalsection][0][]06XY}{[1][0][]source}{}{}{}}
\makeatother
\title{Base change of a K-flat differential graded algebra computes the derived restriction and commutes with the ground-ring derived tensor}
\author{The Stacks Project Authors}
\date{stacks-project:dga.tex:local-snapshot-20261012; GNU FDL 1.2 or later}
\begin{document}
\maketitle
\noindent\textit{Editorial note.} All delivered mathematical text is the author's own, copied from the pinned Stacks Project chapter file sources/stacks-dga/pinned-original-source.tex; no statement or proof step was rewritten, reordered or removed.\par\medskip
\noindent\textit{Editorial note.} The counted proof is the author's complete argument: the canonical comparison map c\_M, the claim for property-(P) modules from which the three parts follow through the derived-tensor definitions, the propagation of isomorphism and K-flatness through sums of shifts, the induction over the split filtration with more-algebra-lemma-K-flat-two-out-of-three, and the filtered-colimit conclusion are all retained verbatim.\par\medskip
\noindent\textit{Editorial note.} The label lemma-base-change-K-flat names the counted statement, and the proof environment immediately following it in the source is delivered verbatim as its proof.\par\medskip
\noindent\textit{Editorial note.} The counted statement is carried by the author's own lemma environment in the chapter source.\par\medskip
\noindent\textit{Allowed standard prerequisites (editorial).}\begin{itemize}
\item Derived tensor products through property (P) resolutions. Same chapter, Lemma lemma-derived-bc; More on Algebra, Section more-algebra-section-derived-tensor-product. Fixes the definition of the derived tensor functors, so that the claim about property-(P) modules converts into the three stated parts by the defining formalism.
\item K-flatness bookkeeping for sums, extensions and colimits. More on Algebra, Lemmas more-algebra-lemma-colimit-K-flat and more-algebra-lemma-K-flat-two-out-of-three. Propagate K-flatness and the isomorphism c from the algebra through sums of shifts, the split filtration extensions, and the final filtered colimit.
\end{itemize}\medskip
\ifdefined\hypertarget\hypertarget{sourcepointer-1}{}\fi
\noindent\textit{Source reference.} Original reference 09LS: sources/stacks-dga/pinned-original-source.tex, source lines 5541--5541. Prerequisite pointer invoked at proof line(s) 5757; it names a standard fact of the same chapter that the counted argument uses, and every step of the counted proof built on it is retained verbatim in the delivered body.\par\medskip
\ifdefined\hypertarget\hypertarget{sourcepointer-2}{}\fi
\noindent\textit{Source reference.} Original lemma 06Y2: sources/stacks-chow/tags.txt, source lines 8871--8871. Prerequisite pointer invoked at proof line(s) 5775; it names a standard fact of a sibling chapter that the counted argument uses, and every step of the counted proof built on it is retained verbatim in the delivered body.\par\medskip
\ifdefined\hypertarget\hypertarget{sourcepointer-3}{}\fi
\noindent\textit{Source reference.} Original lemma 06Y3: sources/stacks-chow/tags.txt, source lines 8872--8872. Prerequisite pointer invoked at proof line(s) 5765; it names a standard fact of a sibling chapter that the counted argument uses, and every step of the counted proof built on it is retained verbatim in the delivered body.\par\medskip
\ifdefined\hypertarget\hypertarget{sourcepointer-4}{}\fi
\noindent\textit{Source reference.} Original section 06XY: sources/stacks-chow/tags.txt, source lines 8867--8867. Prerequisite pointer invoked at proof line(s) 5758; it names a standard fact of a sibling chapter that the counted argument uses, and every step of the counted proof built on it is retained verbatim in the delivered body.\par\medskip
\begin{lemma}
\label{lemma-base-change-K-flat}
Let $R \to R'$ be a ring map. Let $(A, \text{d})$ be a differential
graded $R$-algebra. Let $(A', \text{d})$ be the base change, i.e.,
$A' = A \otimes_R R'$. If $A$ is K-flat as a complex of $R$-modules,
then
\begin{enumerate}
\item $- \otimes_A^\mathbf{L} A' : D(A, \text{d}) \to D(A', \text{d})$
is equal to the right derived functor of
$$
K(A, \text{d}) \longrightarrow K(A', \text{d}),\quad
M \longmapsto M \otimes_R R'
$$
\item the diagram
$$
\xymatrix{
D(A, \text{d}) \ar[r]_{- \otimes_A^\mathbf{L} A'} \ar[d]_{restriction} &
D(A', \text{d}) \ar[d]^{restriction} \\
D(R) \ar[r]^{- \otimes_R^\mathbf{L} R'} & D(R')
}
$$
commutes, and
\item if $M$ is K-flat as a complex of $R$-modules, then the
differential graded $A'$-module $M \otimes_R R'$ represents
$M \otimes_A^\mathbf{L} A'$.
\end{enumerate}
\end{lemma}

\begin{proof}
For any differential graded $A$-module $M$ there is a canonical map
$$
c_M : M \otimes_R R' \longrightarrow M \otimes_A A'
$$
Let $P$ be a differential graded $A$-module with property (P).
We claim that $c_P$ is an isomorphism and that $P$ is K-flat
as a complex of $R$-modules. This will prove all the results
stated in the lemma by formal arguments using the definition
of derived tensor product in Lemma \ref{lemma-derived-bc} and
More on Algebra, Section \ref{more-algebra-section-derived-tensor-product}.

\medskip\noindent
Let $F_\bullet$ be the filtration on $P$ showing that $P$ has property (P).
Note that $c_A$ is an isomorphism and $A$ is K-flat as a complex
of $R$-modules by assumption. Hence the same is true for
direct sums of shifts of $A$ (you can use
More on Algebra, Lemma \ref{more-algebra-lemma-colimit-K-flat}
to deal with direct sums if you like).
Hence this holds for the complexes $F_{p + 1}P/F_pP$.
Since the short exact sequences
$$
0 \to F_pP \to F_{p + 1}P \to F_{p + 1}P/F_pP \to 0
$$
are split exact as sequences of graded modules, we can argue
by induction that $c_{F_pP}$ is an isomorphism for all $p$
and that $F_pP$ is K-flat as a complex of $R$-modules (use
More on Algebra, Lemma \ref{more-algebra-lemma-K-flat-two-out-of-three}).
Finally, using that $P = \colim F_pP$ we conclude that
$c_P$ is an isomorphism and that $P$ is K-flat
as a complex of $R$-modules (use
More on Algebra, Lemma \ref{more-algebra-lemma-colimit-K-flat}).
\end{proof}

\end{document}
